\chapter{Insulin Devices}

\section*{Definition and Overview}
Insulin devices are the tools used to deliver insulin to people with diabetes. They range from simple syringes and insulin pens to insulin pumps and automated systems that combine pumps with continuous glucose monitoring. The choice of device depends on the type of diabetes, the person's lifestyle, preferences, and clinical needs, and on what is funded. Modern devices have made insulin delivery easier, more accurate, and less painful, and the development of hybrid closed-loop systems, sometimes described as artificial pancreas technology, has transformed glucose control for many people. Correct technique, rotation of injection sites, and proper storage of insulin are essential for safe and effective treatment.

\section*{Epidemiology}
Around 1.3 million Australians live with diabetes, and a significant proportion require insulin. Insulin is essential for everyone with type 1 diabetes and for many people with type 2 diabetes as their condition progresses. The availability of subsidised insulin pumps and continuous glucose monitors through the National Diabetes Services Scheme has increased their use, particularly in children and young people with type 1 diabetes, for whom pump therapy is often the preferred option. Pen devices are the most widely used method of insulin delivery worldwide, and their convenience has improved adherence.

\section*{Aetiology and Pathophysiology}
Diabetes results from inadequate insulin action: in type 1 diabetes the pancreatic beta cells are destroyed and insulin is absent, while in type 2 diabetes insulin resistance and progressive beta cell failure lead to relative deficiency. Insulin must be replaced by injection or infusion because it is destroyed in the stomach if taken by mouth. Different insulins mimic normal physiology: basal insulins provide a steady background level, and bolus insulins cover meals. Devices deliver these insulins with varying precision: syringes and pens deliver set doses, while pumps infuse a continuous basal rate with user-activated boluses. Continuous glucose monitoring linked to pumps enables automated adjustment, smoothing glucose levels and reducing dangerous lows.

\section*{Clinical Features}
\begin{itemize}
    \item Syringes: the traditional method, still used where needed
    \item Insulin pens: pre-filled or refillable, convenient and accurate to half units
    \item Insulin pumps: small wearable devices infusing insulin through a cannula, with programmable rates
    \item Continuous glucose monitors (CGMs): sensors worn on the skin reporting glucose continuously
    \item Hybrid closed-loop systems: pumps that automatically adjust insulin using CGM data
    \item Injection site care: rotation of sites to prevent lumps and fat changes
    \item Insulin storage: refrigeration before opening and safe storage in use
    \item Safe disposal of needles and sharps
\end{itemize}

\section*{Differential Diagnosis}
The choice of insulin device is individualised, and various factors are weighed.
\begin{itemize}
    \item \textbf{Type of diabetes and insulin regimen:} multiple daily injections versus pump therapy
    \item \textbf{Age and life stage:} children, pregnancy, and older people have different needs
    \item \textbf{Lifestyle:} work, sport, shift work, and travel
    \item \textbf{Hypoglycaemia awareness and risk:} automated systems reduce lows
    \item \textbf{Hand function and dexterity:} some devices are easier to use
    \item \textbf{Cost and funding:} access through the National Diabetes Services Scheme
    \item \textbf{Preference and training:} success depends on education and support
\end{itemize}

\section*{When to Seek Medical Care}
People with diabetes should have their insulin regimen and devices reviewed regularly by their diabetes team, and should seek review for difficulty controlling glucose, frequent hypoglycaemia, or problems with their device. New symptoms, illness, or planned surgery require adjustment of insulin. Anyone starting insulin needs structured education and follow-up.

\textbf{Call 000 (or local emergency services) for:}
\begin{itemize}
    \item Severe hypoglycaemia: unconsciousness, seizures, or inability to swallow or treat
    \item Signs of diabetic ketoacidosis: high glucose with vomiting, abdominal pain, rapid breathing, or confusion
    \item A device malfunction that cannot be fixed, with rising glucose and ketones
    \item Severe allergic reactions to insulin or the device
\end{itemize}

\section*{Diagnosis and Investigations}
\begin{itemize}
    \item \textbf{Review of glucose levels:} blood glucose and continuous glucose monitoring data
    \item \textbf{HbA1c testing:} the measure of average glucose over three months
    \item \textbf{Assessment of hypoglycaemia:} frequency, severity, and awareness
    \item \textbf{Ketone testing:} when illness or high glucose occurs
    \item \textbf{Review of technique:} injection sites, timing, and device use
    \item \textbf{Education and training:} structured diabetes education for the chosen device
    \item \textbf{Device review:} suitability, funding, and support needs
\end{itemize}

\section*{Management}
\begin{itemize}
    \item \textbf{Insulin pens and syringes:} multiple daily injections with basal and bolus insulins
    \item \textbf{Insulin pumps:} continuous subcutaneous infusion with flexible basal rates and boluses
    \item \textbf{Continuous glucose monitoring:} sensors with alarms for high and low glucose
    \item \textbf{Hybrid closed-loop systems:} automated insulin delivery for suitable people
    \item \textbf{Carbohydrate counting:} matching bolus doses to meals
    \item \textbf{Site rotation:} preventing lipohypertrophy and improving absorption
    \item \textbf{Hypoglycaemia management:} rapid-acting glucose and glucagon for severe events
    \item \textbf{Diabetes education:} ongoing support from credentialled diabetes educators and specialists
    \item \textbf{Access and funding:} applying through the National Diabetes Services Scheme
\end{itemize}

\section*{Complications}
\begin{itemize}
    \item Hypoglycaemia, the most common complication of insulin therapy
    \item Diabetic ketoacidosis from missed insulin or device failure
    \item Lipohypertrophy, lumps of fat at overused injection sites, which alter absorption
    \item Skin infections at injection or infusion sites
    \item Device problems: occlusion, dislodgement, and battery or sensor failures
    \item Cost and the burden of managing technology
    \item Psychological effects of living with diabetes and its devices
\end{itemize}

\section*{Prognosis and Outcomes}
With modern devices, most people with diabetes can achieve glucose levels close to target and dramatically reduce the risk of complications. Continuous glucose monitoring and automated insulin delivery reduce hypoglycaemia and improve time in range, and their benefits are greatest in people with type 1 diabetes, including children. Insulin therapy is lifelong for type 1 diabetes, but the technology makes it easier and safer than ever. For type 2 diabetes, insulin can be introduced simply and effectively when oral medicines no longer control glucose.

\section*{Prevention and Self-Care}
\begin{itemize}
    \item Attend diabetes education to master your device
    \item Rotate injection and infusion sites every time
    \item Check insulin storage and expiry
    \item Monitor glucose as recommended, including during illness
    \item Carry rapid-acting glucose for hypos and know how to treat them
    \item Have a plan for sick days and device failure
    \item Keep spare supplies and know your emergency contacts
    \item Dispose of sharps safely
    \item Attend regular reviews with your diabetes team
    \item Discuss new technology options with your specialist
\end{itemize}

\section*{Sources and Further Reading}
The following reputable sources provide further information for healthcare students and patients seeking reliable, current advice about insulin devices.
\begin{itemize}
    \item Diabetes Australia. \textit{Insulin and delivery devices.} https://www.diabetesaustralia.com.au/
    \item National Diabetes Services Scheme. \textit{Products and funding.} https://www.ndss.com.au/
    \item Australian Diabetes Educators Association. \textit{Diabetes education services.} https://www.adea.com.au/
    \item Healthdirect Australia. \textit{Insulin and diabetes.} https://www.healthdirect.gov.au/
    \item JDRF Australia. \textit{Technology for type 1 diabetes.} https://jdrf.org.au/
    \item Mayo Clinic. \textit{Insulin therapy and devices.} https://www.mayoclinic.org/
\end{itemize}
